% The research corpus (ci/research.json): a history paper in the Chicago
% notes and bibliography style (biblatex-chicago through biber), on
% biblatex-chicago's own test database (notes-test.bib, from TeX Live).
% TeX Live's cms-notes-sample.tex needs Gentium Plus, which TeX Live
% does not ship.
\documentclass[12pt]{article}
\usepackage[T1]{fontenc}
\usepackage{lmodern}
\usepackage[american]{babel}
\usepackage[autostyle]{csquotes}
\usepackage[notes,backend=biber,strict,compresspages]{biblatex-chicago}
\addbibresource{notes-test.bib}
\title{Trade, Letters and Empire}
\author{M. Historian}
\date{January 1, 2026}
\begin{document}
\maketitle
\section{Introduction}
Historians of the Atlantic world have read its letters closely.\autocite{centinel:letters}
The cotton trade shaped it,\autocite[45--47]{cotton:manufacture} as did
the courts of China\autocite{boxer:china} and Peru.\autocite{babb:peru}
\section{Sources}
Aristotle is cited in Greek\autocite{aristotle:metaphy:gr} and in
translation;\autocite{aristotle:metaphy:trans} Chaucer's life records
too.\autocite{chaucer:liferecords} A second citation of a work is
shortened.\autocite{boxer:china} The style's own manual
says so.\autocite[14.23]{chicago:manual}
\section{Conclusion}
\textcite{dunn:revolutions} and \textcite{davenport:attention} close the
argument.
\printbibliography
\end{document}
